\section{Introduction}

Large language model (LLM)-based agents are increasingly used for tasks that go
beyond single-turn question answering: they call external tools and APIs,
retrieve and incorporate documents, read and write persistent memory, and plan
a sequence of actions toward a goal. In multi-agent systems, several such
agents are composed into a single pipeline, for example, a planner that
decomposes the task, a retriever or tool-using agent that gathers external
information, a worker agent that carries out the core task, and a reviewer or
guard agent that checks the result before it is returned to the user.

Although this decomposition improves modularity and often enhances task performance,
it also introduces internal communication channels, such as inter-agent messages, tool-call arguments, memory writes, and logs,
that are not directly visible to the end user.

Existing evaluations of prompt-injection robustness for LLM-integrated systems
typically focus on whether the final, user-facing response leaks sensitive
content or deviates from the intended task \cite{greshake2023indirect,liu2024formalizing}.
In an agentic or multi-agent pipeline, however, a sensitive value can pass
through, and be exposed at several internal points before any final answer is
produced. An attacker who injects malicious content into a retrieved document,
a tool's output, a task input, or a message exchanged between agents may cause
a secret to surface in a tool argument, an inter-agent message, or a memory
entry, even when the system's final response to the user appears unremarkable.
Evaluating only the final output can therefore substantially underestimate how
exposed such a system really is, and may also miss attacks that degrade task
performance or increase operational cost without leaking anything at all.

This project investigates that gap directly. We ask (i) whether an
automated, iterative red-teaming process can surface effective attack
variants that exploit the internal channels of agentic and multi-agent
systems under a small search budget, and (ii) whether final-output-only
evaluation indeed misses leakage and degradation that an internal-channel-aware
evaluation would catch. We adopt an
AutoResearch-style loop, described below: a heuristic search procedure
generates an attack variant, the variant is run against a target agentic
system on a fixed task, the outcome is scored against a chosen metric, and
the variant is kept or discarded before the next is generated. In the reported
evaluation, we operationalize these risks through three experimental goals:
general channel-aware leakage, external final-output leakage, and performance
degradation. The framework can also amplify operational cost through extra
tool calls and token usage; although latency can be recorded as a diagnostic
signal, the reported evaluation focuses on task success, leakage, tool-call
increase, and operational degradation score rather than latency. We study these
goals across four systems---AgentDojo, AutoGen, CrewAI, and LangGraph---configured
with a shared abstract pipeline so that results are comparable across
frameworks. Throughout, the attacker only injects content into untrusted
channels (retrieved documents, tool outputs, task inputs, inter-agent
messages), no real user data is involved and all sensitive information is
represented by synthetic canary secrets in sandboxed tasks.

\subsection{Agentic and Multi-Agent LLM Systems}

An agentic LLM system augments a base language model with the ability to call
external tools or APIs, retrieve and incorporate external documents, read and
write to a memory store, and plan over multiple steps rather than respond in a
single turn. AgentDojo \cite{agentdojo2024} is one such environment: it
provides realistic tool-using tasks together with an established benchmark for
prompt-injection attacks and defenses, and serves in this project as a
single-agent, tool-using reference point.

Multi-agent frameworks compose several agents of this kind into a single
pipeline. AutoGen \cite{autogen2023} structures this composition as a
conversation in which agents exchange messages and may invoke tools or request
human input. CrewAI \cite{crewai} assigns each agent an explicit role, goal,
and set of tools, and lets agents delegate subtasks to one another. LangGraph
\cite{langgraph} represents the pipeline as an explicit graph of nodes and
edges with shared, inspectable state, which makes it convenient to trace
exactly where in the pipeline an attack takes effect. In these multi-agent
frameworks, additional roles such as planner, retriever/tool agent, worker,
and reviewer/guard communicate through messages, shared memory, or delegation
calls that never appear in the final response shown to the user, but that
attacker-controlled input may still be able to influence or read from. The
threat model introduced later in this paper targets exactly these channels.

\subsection{AutoResearch-Style Optimization}

AutoResearch \cite{autoresearch} denotes an iterative research-automation
pattern in which an LLM agent edits an experimental configuration, prompt, or
piece of code, runs an evaluation, scores the outcome against a target metric,
and keeps the change only if it improves that metric, otherwise discarding it
and proposing a different change. The pattern was originally proposed to
automate iterative improvement of training and evaluation pipelines, but the
underlying generate-evaluate-score-keep/discard cycle does not depend on what
is being optimized.

This project adapts that cycle to automated red teaming, replacing the
LLM-agent edit step with a heuristic search procedure. The search selects,
mutates, and recombines variants within a predefined attack space of structured
fields and templates; it does not rely on an LLM to invent the attack space or
families from scratch. The search explores that space, covering the injected
content, its placement within the pipeline, the targeted agent, and its
phrasing, using mutation, crossover, and family-based selection to bias future
variants toward higher-scoring regions of that space. This differs from recent
reinforcement-learning-based approaches to automated prompt injection, which
optimize universal adversarial suffixes through gradient-based or
policy-gradient training over model weights
\cite{learningtoinject2026,verifiableagenttoagent2026}. Because the search
operates on discrete natural-language attack variants rather than on model
parameters, it requires no gradient access and can, in principle, be applied
directly to closed, black-box agentic systems.

\subsection{Contributions}

This project makes the following contributions:
\begin{itemize}
    \item We instantiate an AutoResearch-style red-teaming loop for agentic
    and multi-agent LLM systems, in which a heuristic search procedure
    iteratively generates, evaluates, and selects attack variants against a
    fixed target system and task.
    \item We define three experimental attack goals: general information
    leakage across monitored channels, external leakage into the final
    user-visible output, and performance degradation.
    \item We define an evaluation methodology that separately measures
    final-output leakage and internal-channel leakage (in tool arguments,
    inter-agent messages, memory writes, and logs), enabling a direct test of
    whether output-only evaluation underestimates real exposure.
    \item We apply this methodology to four systems spanning single-agent
    (AgentDojo) and multi-agent (AutoGen, CrewAI, LangGraph) architectures,
    configured with a shared abstract workflow so that leakage and
    degradation behavior can be compared across frameworks.
\end{itemize}
